\documentclass[12pt]{article}
\usepackage[utf8]{inputenc}
\usepackage[T1]{fontenc}
\usepackage{hyperref}
\title{Thakali Language Resources for K-12 Curriculum Development}
\author{Compiled by Claude Sonnet 5 (claude-sonnet-5) for Andr\'as Kornai}
\date{}
\begin{document}
\maketitle

\section{Introduction}
Thakali (ISO 639-3 ths) is a Sino-Tibetan language of Nepal, spoken mainly in Mustang and Myagdi. Speaker counts differ by source: the 2011 census figure given is 5,242, and Wikipedia lists about 4,200. This report is based on a small number of searches. Very little is available online.

\section{Primary Educational Resources}
A 2010 SEAMEO presentation says Nepal's Curriculum Development Centre has produced reference materials in Thakali. It does not list Thakali among the grade 1--5 mother-tongue books. The information is old and should be rechecked. \url{https://www.seameo.org/LanguageMDGConference2010/doc/presentations/day2/IndreshThakur-sn.pdf}

\section{Linguistic and Reference Materials}
\begin{itemize}
\item Mark Turin's research notes mention about 2,000 collected words and a planned Thakali--Nepali dictionary. I found no evidence it was published. \url{https://www.digitalhimalaya.com/projectteam/turin/abstracts/qcr.html}
\item Wikipedia entry on the language. \url{https://en.wikipedia.org/wiki/Thakali_language}
\item Omniglot overview. It says Thakali is taught in primary schools, which I could not confirm. \url{https://www.omniglot.com/writing/thakali.htm}
\end{itemize}

\section{Religious and Community Resources}
\begin{itemize}
\item Indigenous Voice page on the Thakali people. \url{https://english.indigenousvoice.com/thakali}
\item Turin, ``Mind Your Langs''. It discusses a Mustang school with a mother-tongue emphasis that teaches Tibetan rather than Thakali. \url{https://markturin.sites.olt.ubc.ca/files/2018/07/Turin_Mind_Your_Langs.pdf}
\end{itemize}

\section{Digital and Audio Resources}
None were found in the search results.

\section{Curriculum Development Context}
There is almost no published teaching material. A curriculum would start from documentation, a dictionary and an orthography agreed with the community. Current school status should be checked with the Curriculum Development Centre and Thakali community organisations.
\end{document}